% ADR-001: Microservices Architecture
\subsection{ADR-001: Adoption of Microservices Architecture}

\begin{tabular}{p{3cm}p{10cm}}
\textbf{Status} & Accepted \\
\textbf{Date} & 2026-03-15 \\
\textbf{Decision} & Adopt a microservices architecture over a monolithic architecture. \\
\end{tabular}

\textbf{Context}: The Eventify platform comprises distinct domains: user management, event management, real-time chat, and notifications. These domains have different scaling characteristics. Event browsing may experience high read traffic, while chat requires persistent WebSocket connections. A monolithic application would force all domains to scale together, wasting resources.

\textbf{Decision}: Decompose the system into five independently deployable microservices: API Gateway, User Service, Event Service, Chat Service, and Notification Service. Each service owns its data and communicates via REST (synchronous) or RabbitMQ (asynchronous).

\textbf{Consequences}:
\begin{itemize}
    \item \textit{Positive}: Independent scaling, independent deployment, technology flexibility per service, fault isolation.
    \item \textit{Negative}: Increased operational complexity, network latency between services, eventual consistency challenges, need for distributed tracing.
\end{itemize}

\textbf{Alternatives Considered}:
\begin{itemize}
    \item Monolithic architecture --- simpler to develop initially but does not meet scalability and independent deployment requirements.
    \item Service-oriented architecture (SOA) --- heavier governance and shared data layer; microservices provide better autonomy.
\end{itemize}
